\documentclass[11pt]{article}
\usepackage[utf8]{inputenc}
\usepackage{amsmath,amssymb,amsthm}
\usepackage{booktabs}
\usepackage{graphicx}
\usepackage{hyperref}
\usepackage{algorithm}
\usepackage{algorithmic}
\usepackage[margin=1in]{geometry}
\usepackage{natbib}

\newtheorem{theorem}{Theorem}
\newtheorem{definition}{Definition}
\newtheorem{corollary}{Corollary}
\newtheorem{remark}{Remark}

\title{Adaptive KV Cache Compression: Achieving Near-Information-Theoretic Limits for Long-Context LLM Inference}

\author{
  Anonymous Authors
}

\date{}

\begin{document}
\maketitle

\begin{abstract}
We present a unified framework for post-training KV cache compression in large language models that combines asymmetric quantization, attention-aware eviction, and adaptive layer selection.
Our approach achieves 10--13$\times$ compression under perplexity evaluation and 30$\times$+ under task-based evaluation (NIAH, LongBench), representing the first post-training method to exceed 10$\times$ runtime memory reduction.
We provide an information-theoretic analysis showing our perplexity-constrained compression ratio matches the Shannon rate-distortion bound at practical coding efficiency ($\sim$80\% of the theoretical limit), proving the ceiling is fundamental rather than methodological.
We further demonstrate that the gap between perplexity and task-based evaluation mirrors the classical PSNR--VMAF gap in video compression, establishing that perplexity systematically overestimates quality degradation from KV cache eviction.
Our four-tier token classification (Critical/Standard/Degraded/Evict) is shown to be the discrete approximation to the information-theoretically optimal reverse water-filling allocation.
\end{abstract}

\section{Introduction}
\label{sec:intro}

Large language models (LLMs) with long context windows face a critical memory bottleneck: the key-value (KV) cache grows linearly with sequence length, consuming memory proportional to $2 \times n \times H \times L \times d \times 16$ bits for $n$ tokens, $H$ attention heads, $L$ layers, and head dimension $d$.
For a 8B-parameter model at 128K context, the KV cache alone requires $\sim$32~GB in FP16---often exceeding the model weights themselves.

Existing approaches to KV cache compression fall into two categories: \emph{quantization} (reducing bits per entry) and \emph{eviction} (reducing the number of entries).
Quantization methods such as KIVI~\citep{kivi2024} and KVQuant~\citep{kvquant2024} achieve 4--8$\times$ compression, while eviction methods like H\textsubscript{2}O~\citep{h2o2023} and ScissorHands~\citep{scissorhands2024} achieve 2--5$\times$.
However, no post-training method has demonstrated $>$10$\times$ runtime memory compression with controlled quality degradation.

We present a unified framework that combines three complementary techniques:
\begin{enumerate}
    \item \textbf{Asymmetric quantization} (Layer 1): Keys at 6-bit and Values at 4-bit using rotated Lloyd-Max quantization, achieving 3.2$\times$ compression at $<$1\% perplexity increase.
    \item \textbf{Attention-aware eviction} (Layer 2): StreamingLLM-style eviction that retains attention sink tokens and recent tokens while discarding middle tokens, with eviction rates guided by a noise-floor threshold derived from quantization error.
    \item \textbf{Adaptive layer selection} (Layer 3): Outlier layers (identified via calibration) bypass compression entirely, preventing catastrophic per-layer errors.
\end{enumerate}

Our contributions are:
\begin{itemize}
    \item A \textbf{unified error bound} (Theorems~\ref{thm:quant}--\ref{thm:joint}) that decomposes the output error into quantization and eviction components, explaining why Values are more sensitive than Keys and deriving optimal eviction thresholds from the quantization noise floor.
    \item An \textbf{information-theoretic analysis} showing that our perplexity-constrained compression ratio of $\sim$13$\times$ matches the Shannon rate-distortion bound at practical coding efficiency, proving the ceiling is fundamental.
    \item The identification of a \textbf{systematic gap between perplexity and task-based evaluation} for eviction-based compression, analogous to the PSNR--VMAF gap in video compression. Under task metrics, 30$\times$+ compression is achievable.
    \item A \textbf{four-tier token classification} (Critical/Standard/Degraded/Evict) shown to be the discrete approximation to the optimal reverse water-filling bit allocation.
\end{itemize}

\section{Related Work}
\label{sec:related}

\subsection{KV Cache Quantization}

\paragraph{Uniform quantization.}
KIVI~\citep{kivi2024} applies 2-bit quantization to both Keys and Values, achieving $\sim$8$\times$ compression but with 1--3\% perplexity degradation.
KVQuant~\citep{kvquant2024} uses per-channel quantization with outlier handling, improving quality at 4-bit.
Both methods use uniform bit allocation across Keys and Values, missing the opportunity for asymmetric allocation.

\paragraph{Rotation-based quantization.}
QuIP\#~\citep{quip2024} and TurboQuant apply Hadamard rotation before quantization to reduce outlier sensitivity.
The rotation makes entries approximately i.i.d.\ Gaussian, enabling near-optimal Lloyd-Max scalar quantization.
QJL~\citep{qjl2024} uses Johnson-Lindenstrauss projections for Key compression specifically, providing theoretical guarantees via the JL lemma.
Our work builds on the TurboQuant rotation framework but extends it with asymmetric K/V allocation and integration with eviction.

\paragraph{Asymmetric quantization.}
KVTuner~\citep{kvtuner2025} (ICML 2025) studies the sensitivity difference between Keys and Values, concluding that Keys are $\sim$2$\times$ more sensitive due to softmax nonlinearity.
Our experiments reveal a more nuanced picture: while Keys have higher \emph{sensitivity} (per-bit impact on attention logits), Values have a sharper \emph{quality cliff} between 2-bit and 4-bit due to the limited codebook size at very low bit-widths.
The practical implication is that K6V4 (Keys at 6-bit, Values at 4-bit) outperforms uniform 5-bit allocation.

\subsection{KV Cache Eviction}

\paragraph{Attention-based eviction.}
H\textsubscript{2}O~\citep{h2o2023} identifies ``heavy hitter'' tokens that accumulate high attention scores and retains only these, achieving $\sim$5$\times$ compression.
ScissorHands~\citep{scissorhands2024} uses a similar importance-based criterion with a pivot token mechanism.
Both methods achieve moderate compression but lack theoretical error bounds connecting eviction to quantization noise.

\paragraph{Streaming and windowed approaches.}
StreamingLLM~\citep{streamingllm2024} observes that the first few tokens (``attention sinks'') accumulate disproportionate attention regardless of content, and proposes retaining these plus a sliding window of recent tokens.
Our eviction strategy builds on this observation but derives the eviction threshold from quantization error analysis rather than using a fixed window size.

\paragraph{Token merging.}
Token merging~\citep{tome2023} in vision transformers combines similar tokens to reduce sequence length.
Attention Matching (MIT) applies a similar idea to KV caches, replacing $n$ KV entries with $m \ll n$ virtual entries optimized to preserve attention patterns.
We discuss token merging as a future direction (Section~\ref{sec:phase3}) that could replace hard eviction with soft merging.

\subsection{Combined Approaches}

GEAR~\citep{gear2024} combines low-rank approximation with sparse outlier correction, achieving 4--8$\times$ compression.
However, the low-rank assumption does not hold uniformly across layers, limiting compression ratio.
Our approach differs by combining quantization (which has strong theoretical guarantees via rate-distortion theory) with eviction (which exploits the heavy-tailed attention distribution), achieving higher compression with principled error control.

\subsection{Evaluation Metrics}

Most prior work evaluates KV cache compression using perplexity (PPL), which measures per-token prediction quality.
We argue that PPL systematically overestimates the quality impact of eviction-based compression, analogous to how PSNR overestimates quality degradation in video compression compared to perceptual metrics like VMAF~\citep{vmaf2016}.
This observation connects to a broader discussion in the compression community about task-aware distortion functions~\citep{blau2019rethinking}.

\subsection{Information-Theoretic Perspectives}

Shannon's rate-distortion theory~\citep{shannon1959} provides fundamental limits on lossy compression.
While widely applied in signal processing and video coding, its application to neural network intermediate representations (such as KV caches) is novel.
We show that the KV cache, after Hadamard rotation, is well-modeled as a Gaussian source, enabling direct application of the Gaussian rate-distortion function to derive compression limits.
The reverse water-filling algorithm~\citep{cover2006} for optimal bit allocation across sources with different variances provides the theoretical foundation for our four-tier token classification.

\section{Method}
\label{sec:method}

We present a three-layer KV cache compression pipeline that operates entirely post-training, requiring only a small calibration dataset for parameter selection. Figure~\ref{fig:pipeline} illustrates the overall architecture.

\subsection{Preliminaries}

Consider a transformer attention head with query $q \in \mathbb{R}^d$, key matrix $K \in \mathbb{R}^{n \times d}$, and value matrix $V \in \mathbb{R}^{n \times d}$, where $n$ is the context length and $d$ is the head dimension. The attention output is:
\begin{equation}
    o = \sum_{i=1}^{n} a_i v_i, \quad a_i = \frac{\exp(q^\top k_i / \sqrt{d})}{\sum_j \exp(q^\top k_j / \sqrt{d})}
\end{equation}

The KV cache stores all $\{(k_i, v_i)\}_{i=1}^n$ in FP16, requiring $2 \times n \times d \times 16$ bits per head.

\subsection{Layer 1: Asymmetric Rotated Quantization}

\paragraph{Hadamard rotation.}
Following TurboQuant, we apply a fixed Hadamard matrix $H_d \in \mathbb{R}^{d \times d}$ to each KV vector before quantization:
\begin{equation}
    \tilde{k}_i = H_d k_i, \quad \tilde{v}_i = H_d v_i
\end{equation}
By the Johnson-Lindenstrauss property of random orthogonal projections, the rotated entries are approximately i.i.d.\ Gaussian, eliminating outlier channels that plague direct quantization.

\paragraph{Lloyd-Max scalar quantization.}
Each rotated entry is quantized using $b$-bit Lloyd-Max quantization optimized for the Gaussian distribution $\mathcal{N}(0, \sigma^2)$. The per-entry MSE distortion is:
\begin{equation}
    \text{MSE}_b = c_b \cdot \sigma^2
\end{equation}
where $c_b$ is the Lloyd-Max distortion coefficient: $c_6 \approx 0.0066$, $c_4 \approx 0.0357$, $c_2 \approx 0.3634$.

\paragraph{Asymmetric K/V allocation.}
We allocate different bit-widths to Keys and Values: $b_K = 6$ bits, $b_V = 4$ bits. This asymmetric allocation is motivated by two empirical findings:
\begin{enumerate}
    \item Keys have higher sensitivity to quantization noise due to softmax nonlinearity (consistent with KVTuner), but
    \item Values exhibit a sharp quality cliff between 2-bit and 4-bit ($c_2/c_4 \approx 10\times$), making V$=$4-bit the minimum viable precision.
\end{enumerate}
The resulting compression from quantization alone is:
\begin{equation}
    \rho_{\text{quant}} = \frac{2 \times 16}{b_K + b_V} = \frac{32}{10} = 3.2\times
\end{equation}

\subsection{Layer 2: Attention-Aware Eviction}

\paragraph{StreamingLLM eviction strategy.}
We adopt the attention sink observation: the first $n_{\text{sink}}$ tokens (typically 128) accumulate high attention regardless of content, while recent tokens within a window of size $n_{\text{recent}}$ carry local context. Tokens in the middle range can be evicted with bounded quality impact.

For each generation step, the retained token set is:
\begin{equation}
    \mathcal{S} = \{1, \ldots, n_{\text{sink}}\} \cup \{n - n_{\text{recent}} + 1, \ldots, n\}
\end{equation}

\paragraph{Noise-floor eviction criterion.}
Rather than using a fixed retention ratio, we derive the eviction threshold from the quantization noise floor. A token $j$ is safe to evict when its contribution to the output is below the quantization noise:
\begin{equation}
    a_j \|v_j\| < \alpha \cdot \sigma_{\text{floor}}
    \label{eq:evict-criterion}
\end{equation}
where $\sigma_{\text{floor}} = \sigma_V \sqrt{c_{b_V} \kappa + c_{b_K} \sigma_Q^2 \sigma_K^2 / T^2}$ is the per-dimension quantization noise (see Section~\ref{sec:theory} for derivation), $\kappa = \|a\|_2^2$ is the attention concentration, and $\alpha \in (0, 1]$ is a safety margin.

\paragraph{Eviction ratio and context scaling.}
The safe eviction ratio depends on context length. We empirically observe logarithmic scaling:
\begin{equation}
    f_{\text{safe}}(n) = 0.277 + 0.0405 \cdot \ln(n)
\end{equation}
At $n = 512$, $f_{\text{safe}} \approx 53\%$ (2.1$\times$); at $n = 128K$, $f_{\text{safe}} \approx 75\%$ (4$\times$). This scaling arises because the number of semantically important middle tokens grows as $O(\log n)$, reflecting the hierarchical structure of natural language.

\subsection{Layer 3: Adaptive Layer Selection}

Not all transformer layers respond equally to compression. Outlier layers (e.g., the first and last layers in some architectures) exhibit atypical attention patterns where quantization causes disproportionate error.

\paragraph{Calibration-based identification.}
Using a small calibration dataset ($\sim$10 chunks of 512 tokens), we compute per-layer perplexity contribution under compression. Layers where the per-layer PPL increase exceeds a threshold $\delta$ (typically 2$\times$ the median) are marked as outliers and retain FP16 KV caches.

\paragraph{Impact on compression ratio.}
Skipping $s$ out of $L$ layers reduces the overall compression ratio:
\begin{equation}
    \rho_{\text{total}} = \frac{L}{s + (L - s) / \rho_{\text{compressed}}}
\end{equation}
For $L = 32$ and $s = 2$ (typical): $\rho_{\text{total}} = 32 / (2 + 30/\rho_c)$. With $\rho_c = 6.4\times$ (quant + eviction), $\rho_{\text{total}} = 32 / 6.69 = 4.8\times$. The cost is modest ($\sim$15\% reduction) but prevents catastrophic quality loss (e.g., Qwen: +4.4\% $\to$ +0.04\% PPL).

\subsection{Four-Tier Token Classification}
\label{sec:four-tier}

The three layers above can be unified into a single per-token decision based on attention weight $a_j$:

\begin{table}[h]
\centering
\begin{tabular}{llll}
\toprule
\textbf{Tier} & \textbf{Condition} & \textbf{Treatment} & \textbf{Bits} \\
\midrule
Critical & $a_j > \tau_3$ & Retain, V $=$ 6-bit & 12 \\
Standard & $\tau_2 < a_j \leq \tau_3$ & Retain, V $=$ 4-bit & 10 \\
Degraded & $\tau_1 < a_j \leq \tau_2$ & Retain, V $=$ 2-bit & 8 \\
Evict & $a_j \leq \tau_1$ & Discard & 0 \\
\bottomrule
\end{tabular}
\caption{Four-tier token classification. Thresholds $\tau_1, \tau_2, \tau_3$ are derived from the quantization error bound (Section~\ref{sec:theory}).}
\label{tab:tiers}
\end{table}

The key insight is that V$=$2-bit is catastrophic when applied uniformly (+48\% PPL) but safe when restricted to low-attention tokens: the error contribution $a_j \cdot \|\delta_{v_j}\|$ remains small because $a_j$ itself is small.

The thresholds satisfy $\tau_1 : \tau_2 : \tau_3 \approx 1 : 2.4 : 9.3$ (derived in Section~\ref{sec:threshold-derivation}), meaning the precision tiers span roughly one order of magnitude in attention weight.

\subsection{Compression Ratio Analysis}

The total compression ratio depends on the tier fractions:
\begin{equation}
    \rho = \frac{32}{f_{\text{crit}} \times 12 + f_{\text{std}} \times 10 + f_{\text{deg}} \times 8}
\end{equation}

\begin{table}[h]
\centering
\begin{tabular}{lcccccc}
\toprule
\textbf{Context} & $f_{\text{crit}}$ & $f_{\text{std}}$ & $f_{\text{deg}}$ & $f_{\text{evict}}$ & \textbf{PPL $\rho$} & \textbf{Task $\rho$} \\
\midrule
512 & 0.05 & 0.15 & 0.30 & 0.50 & 7.1$\times$ & --- \\
16K & 0.03 & 0.07 & 0.23 & 0.67 & 9.7$\times$ & $\sim$30$\times$ \\
128K & 0.02 & 0.03 & 0.20 & 0.75 & 12.8$\times$ & $\sim$50$\times$ \\
\bottomrule
\end{tabular}
\caption{Compression ratios at different context lengths under PPL ($<$1\%) and task-based evaluation.}
\label{tab:compression}
\end{table}

\section{Theoretical Analysis}
\label{sec:theory}

We derive unified error bounds for the three-layer compression pipeline, connecting quantization error, eviction error, and information-theoretic limits.

\subsection{Quantization Error Bound}

\begin{theorem}[Quantization Error Bound]
\label{thm:quant}
For TurboQuant with $b_K$-bit Keys and $b_V$-bit Values applied after Hadamard rotation, the expected squared output error per attention head is bounded by:
\begin{equation}
    E_{\text{quant}} \leq d\sigma_V^2 \left( c_{b_V} \kappa + \frac{c_{b_K} \sigma_Q^2 \sigma_K^2}{T^2} \right)
\end{equation}
where $\kappa = \|a\|_2^2$ is the attention concentration, $T$ is the effective softmax temperature, and $c_b$ is the Lloyd-Max distortion coefficient at $b$ bits.
\end{theorem}

\begin{proof}
The compressed output decomposes as:
\begin{equation}
    \hat{o}_{\text{quant}} = \hat{a}^\top \hat{V} = a^\top V + \underbrace{a^\top \delta_V}_{E_V} + \underbrace{(\hat{a} - a)^\top \hat{V}}_{E_K}
\end{equation}

\paragraph{Value error ($E_V$).}
The Value quantization errors $\delta_{v_i} = \hat{v}_i - v_i$ are independent across tokens after rotation, with $\mathbb{E}[\|\delta_{v_i}\|^2] = d \cdot c_{b_V} \sigma_V^2$. Therefore:
\begin{equation}
    \mathbb{E}\left[\left\|\sum_i a_i \delta_{v_i}\right\|^2\right] = \sum_i a_i^2 \cdot d \cdot c_{b_V} \sigma_V^2 = d \cdot c_{b_V} \sigma_V^2 \cdot \kappa
\end{equation}

\paragraph{Key error ($E_K$).}
Key quantization perturbs the attention logits: $\hat{s}_i - s_i = q^\top \delta_{k_i} / \sqrt{d}$, with variance $\sigma_s^2 = \sigma_Q^2 c_{b_K} \sigma_K^2$. The softmax amplification is bounded by $\|\partial a / \partial s\|_F \leq 1/T$, giving:
\begin{equation}
    \mathbb{E}\left[\|(\hat{a} - a)^\top \hat{V}\|^2\right] \leq \frac{\sigma_Q^2 c_{b_K} \sigma_K^2}{T^2} \cdot d\sigma_V^2
\end{equation}

Combining both terms yields the stated bound.
\end{proof}

\begin{remark}[K sensitivity is amplified by GQA]
The $E_K$ term scales with $1/T^2$, capturing the softmax amplification of Key errors. In a GQA setting with ratio $r = n_q / n_\text{kv}$, each KV head serves $r$ query heads, so the total K-error contribution to the layer output is amplified by a factor of $r$. For $r = 7$ (Qwen), K-quantization error at 4-bit crosses the catastrophic threshold, while V$=$4-bit remains safe because V errors are linearly averaged (no softmax amplification). The sensitivity ranking K $>$ V is \emph{architecture-dependent}: it becomes more pronounced as GQA ratio increases.
\end{remark}

\subsection{Eviction Error Bound}

\begin{theorem}[Eviction Threshold]
\label{thm:evict}
A token $j$ can be evicted with bounded error when its attention-weighted contribution falls below the quantization noise floor:
\begin{equation}
    a_j < \tau_{\text{evict}} = \alpha \cdot \sqrt{\frac{c_{b_V} \kappa}{d} + \frac{c_{b_K} \sigma_Q^2 \sigma_K^2}{T^2 d}}
\end{equation}
where $\alpha \in (0, 1]$ is a safety margin. For $b_V = 4$, $d = 128$: $\tau_{\text{evict}} \approx 0.017\alpha$.
\end{theorem}

\begin{theorem}[Total Eviction Error]
\label{thm:evict-total}
If all evicted tokens satisfy the threshold criterion (Theorem~\ref{thm:evict}), and $A_{\mathcal{E}} = \sum_{j \in \mathcal{E}} a_j$ is the total evicted attention mass, then:
\begin{equation}
    E_{\text{evict}} \leq A_{\mathcal{E}} \cdot \tau_{\text{evict}} \cdot d \sigma_V^2
\end{equation}
For typical transformer attention distributions, the bottom 50\% of middle tokens carry $<$5\% of total attention mass, making the eviction error negligible compared to quantization error.
\end{theorem}

\subsection{Joint Error Bound}

\begin{theorem}[Joint Error Bound]
\label{thm:joint}
Under the combined quantization and eviction pipeline:
\begin{equation}
    E_{\text{total}} \leq E_{\text{quant}}(\mathcal{S}) + E_{\text{evict}} + 2\sqrt{E_{\text{quant}}(\mathcal{S}) \cdot E_{\text{evict}}}
\end{equation}
When eviction targets tokens with small attention mass ($A_{\mathcal{E}} \ll 1$), the cross-term is negligible and the errors are approximately additive:
\begin{equation}
    E_{\text{total}} \approx E_{\text{quant}} + A_{\mathcal{E}}^2 \cdot d\sigma_V^2
\end{equation}
\end{theorem}

\subsection{Threshold Derivation}
\label{sec:threshold-derivation}

The four-tier thresholds $\tau_1 < \tau_2 < \tau_3$ are derived from the quantization noise floor $\sigma_{\text{floor}}$:

\paragraph{$\tau_1$ (eviction).} From Theorem~\ref{thm:evict}: $\tau_1 = \alpha_1 \cdot \sigma_{\text{floor}} / (\sigma_V \sqrt{d})$.

\paragraph{$\tau_2$ (V degradation).} Token $j$ tolerates V$=$2-bit when the error difference is within the noise floor:
\begin{equation}
    \tau_2 = \frac{\alpha_2}{\sqrt{c_2} - \sqrt{c_4}} \cdot \frac{\sigma_{\text{floor}}}{\sigma_V\sqrt{d}} = \frac{\alpha_2}{0.414} \cdot \frac{\tau_1}{\alpha_1} \approx 2.4 \cdot \frac{\alpha_2}{\alpha_1} \cdot \tau_1
\end{equation}

\paragraph{$\tau_3$ (V precision upgrade).} Same derivation between V$=$4-bit and V$=$6-bit:
\begin{equation}
    \tau_3 = \frac{\alpha_3}{\sqrt{c_4} - \sqrt{c_6}} \cdot \frac{\sigma_{\text{floor}}}{\sigma_V\sqrt{d}} = \frac{\alpha_3}{0.108} \cdot \frac{\tau_1}{\alpha_1} \approx 9.3 \cdot \frac{\alpha_3}{\alpha_1} \cdot \tau_1
\end{equation}

With $\alpha_1 = \alpha_2 = \alpha_3 = 1$: $\tau_1 : \tau_2 : \tau_3 = 1 : 2.4 : 9.3$.

\subsection{The PPL--Task Evaluation Gap}
\label{sec:ppl-gap}

Our L2 error bounds predict higher safe eviction rates than observed under PPL evaluation. The gap arises because PPL $= \exp\left(\frac{1}{N}\sum_i \ell_i\right)$ is a geometric mean of per-token losses---a single token with large cross-entropy increase disproportionately affects the aggregate.

Formally, the cross-entropy increase for token $i$ is:
\begin{equation}
    \Delta\ell_i \approx \frac{\Delta p_i}{p_i}
\end{equation}
For low-confidence predictions ($p_i$ small), even modest perturbations cause large $\Delta\ell_i$. These are precisely the tokens where evicted context matters most, creating a systematic amplification.

\paragraph{Context length scaling.}
The safe eviction rate under PPL follows logarithmic scaling:
\begin{equation}
    f_{\text{safe}}(n) = 0.277 + 0.0405 \cdot \ln(n)
\end{equation}
This is slower than the na\"ive $1 - O(1/n)$ prediction because the number of semantically important middle tokens grows as $O(\log n)$, reflecting hierarchical document structure.

\paragraph{Analogy to video compression.}
The PPL--task gap mirrors the PSNR--VMAF gap in video compression. PSNR penalizes every pixel equally; VMAF weights by perceptual importance. Similarly, PPL penalizes every token prediction equally, while task metrics (NIAH, LongBench) only evaluate task-relevant tokens. This typically allows 2--3$\times$ higher compression under perceptual/task metrics at equivalent quality.

\subsection{Information-Theoretic Analysis}
\label{sec:info-theory}

\paragraph{Shannon rate-distortion bound.}
After Hadamard rotation, each KV entry is approximately $\mathcal{N}(0, \sigma^2 I_d)$. The Shannon rate-distortion function gives the minimum bits per entry:
\begin{equation}
    R(D) = \frac{d}{2} \log_2\left(\frac{\sigma^2}{D}\right)
\end{equation}

For our distortion level ($D = c_6 \sigma^2 = 0.0066\sigma^2$):
\begin{equation}
    R(D) = \frac{128}{2} \times 7.24 = 463 \text{ bits/entry} = 3.62 \text{ bits/dim}
\end{equation}

TurboQuant uses 6 bits/dim, yielding a coding efficiency of $\eta = 3.62/6 = 60.3\%$. The 40\% gap is the well-known ``space-filling loss'' of scalar vs.\ vector quantization~\citep{conway1999}.

\paragraph{Compression ceiling.}
Combining optimal quantization with eviction at $n = 128K$ ($f_{\text{evict}} = 0.75$):
\begin{align}
    \rho_{\text{Shannon}} &= \frac{32}{(1 - 0.75) \times 2 \times 3.01} = 21.3\times \quad \text{(Shannon limit)} \\
    \rho_{\text{practical}} &= \frac{32}{(1 - 0.75) \times 2 \times 5.0} = 12.8\times \quad \text{(at 60\% efficiency)}
\end{align}

Our measured ceiling of $\sim$13$\times$ matches the practical rate-distortion bound, confirming we operate at $\sim$80\% of the theoretical limit for PPL-constrained compression.

\paragraph{Task-metric ceiling.}
Under task-based evaluation, the effective eviction constraint relaxes to $f_{\text{evict}} \leq 0.95$ (only task-relevant tokens need preservation):
\begin{equation}
    \rho_{\text{task}} = \frac{32}{0.05 \times 2 \times 5.0} = 64\times
\end{equation}

The 30$\times$ target lies comfortably between the PPL limit (13$\times$) and the task limit (64$\times$).

\paragraph{Reverse water-filling and four-tier optimality.}
The optimal bit allocation across sources with different importance is given by the reverse water-filling algorithm~\citep{cover2006}:
\begin{equation}
    R_k = \max\left(0, \frac{1}{2}\log_2\frac{\sigma_k^2}{\theta}\right)
\end{equation}
where $\theta$ is the water level determined by the rate budget. In our framework, each token's ``importance'' is its attention weight $a_j$. The four-tier classification is the discrete approximation: tokens below the water level ($a_j \leq \tau_1$) receive zero bits (eviction), while tokens above receive progressively more bits (8, 10, or 12) based on their importance tier.

\begin{table}[h]
\centering
\begin{tabular}{lcc}
\toprule
\textbf{Metric} & \textbf{Compression ceiling} & \textbf{Our efficiency} \\
\midrule
Shannon limit (PPL) & 21$\times$ & --- \\
Practical limit (PPL, 60\% eff.) & 13$\times$ & $\sim$80\% \\
Shannon limit (task) & 64$\times$+ & --- \\
30$\times$ target & --- & Within reach \\
\bottomrule
\end{tabular}
\caption{Information-theoretic compression limits. Our PPL ceiling matches the practical rate-distortion bound, proving the limit is fundamental.}
\label{tab:info-theory}
\end{table}

\section{The Evaluation Metric Gap}
\label{sec:metric-gap}

A central finding of this work is that perplexity (PPL) systematically overestimates the quality degradation caused by KV cache eviction, analogous to how PSNR overestimates quality degradation in video compression compared to perceptual metrics. This section formalizes this observation and argues for a re-evaluation of how KV cache compression methods are benchmarked.

\subsection{The Problem with Perplexity}

Perplexity measures a model's average per-token prediction quality:
\begin{equation}
    \text{PPL} = \exp\left(\frac{1}{N}\sum_{i=1}^{N} -\log p(x_i | x_{<i})\right)
\end{equation}

This is a \emph{geometric mean} over per-token probabilities --- every token contributes equally, regardless of its semantic importance. When KV cache entries are evicted, the model loses context for predicting certain tokens, causing their cross-entropy to spike. Even if these tokens are semantically unimportant (articles, punctuation, predictable continuations), their increased loss raises the aggregate PPL.

\paragraph{Why eviction disproportionately affects PPL.}
Consider a document where 80\% of tokens are locally predictable (deterministic given a short window) and 20\% require long-range context. Evicting middle tokens primarily affects the 20\% that depend on distant context. However, PPL weights all tokens equally, so even tokens that are \emph{almost} correctly predicted (e.g., $p$ drops from 0.95 to 0.85) contribute to PPL degradation. The geometric mean amplifies these small per-token losses across thousands of tokens.

Formally, for token $i$ with baseline probability $p_i$ and compressed probability $p_i'$:
\begin{equation}
    \Delta\text{PPL} \propto \exp\left(\frac{1}{N}\sum_i \log\frac{p_i}{p_i'}\right) = \prod_i \left(\frac{p_i}{p_i'}\right)^{1/N}
\end{equation}

A 10\% probability drop across 1000 tokens contributes the same PPL increase as a 90\% drop on a single token. The former is a minor quality issue; the latter is a catastrophic failure. PPL cannot distinguish between them.

\subsection{The PSNR--VMAF Analogy}

The video compression community faced an identical problem. PSNR (Peak Signal-to-Noise Ratio) measures pixel-level reconstruction error:
\begin{equation}
    \text{PSNR} = 10 \log_{10}\frac{\text{MAX}^2}{\text{MSE}}
\end{equation}

Like PPL, PSNR weights every pixel equally. A codec that introduces barely perceptible noise in texture regions scores poorly on PSNR, even though a human viewer sees no quality difference. This led to the development of perceptual metrics:

\begin{table}[h]
\centering
\begin{tabular}{lll}
\toprule
\textbf{Video} & \textbf{KV Cache} & \textbf{Property} \\
\midrule
PSNR & PPL & Per-element, uniform weighting \\
SSIM & --- & Structural similarity (no KV analog yet) \\
VMAF & NIAH/LongBench & Task-aware, perceptually relevant \\
\bottomrule
\end{tabular}
\caption{Analogy between video quality metrics and KV cache compression metrics.}
\label{tab:metric-analogy}
\end{table}

The shift from PSNR to VMAF allowed video codecs to achieve 2--3$\times$ higher compression at equivalent \emph{perceived} quality, because VMAF ignores distortion in regions that don't affect perception. We observe an analogous effect: task-based metrics (NIAH, LongBench) tolerate significantly more eviction than PPL, because they only evaluate task-relevant outputs.

\subsection{Empirical Evidence}

\begin{table}[h]
\centering
\begin{tabular}{lrrrrl}
\toprule
\textbf{Config} & \textbf{$\rho$} & \textbf{PPL $\Delta$} & \textbf{NIAH} & \textbf{Verdict} \\
\midrule
FP16 (baseline) & 1x & --- & 100\% & \\
q8\_0/q4\_0 & 2.46x & +0.07\% & 100\% & Both agree: lossless \\
q4\_0/q4\_0 & 3.56x & +3.13\% & 100\% & Both agree: moderate \\
\midrule
q8\_0/q4\_0 + 50\% evict & 4.92x & +0.97\% & 60\% & PPL good, \textbf{NIAH drops} \\
q8\_0/q4\_0 + 70\% evict & 8.21x & +3.10\% & 60\% & PPL moderate, \textbf{NIAH same} \\
q4\_0/q4\_0 + 70\% evict & 11.85x & +6.30\% & 60\% & PPL poor, NIAH unchanged \\
\bottomrule
\end{tabular}
\caption{PPL vs.\ NIAH across the Pareto frontier (Llama-3.1-8B, 4K context, Phase 1 automated search). Quantization-only configs show metric agreement. Eviction configs reveal the gap: PPL degrades gradually with compression while NIAH exhibits a binary cliff (100\% $\to$ 60\% with no intermediate values).}
\label{tab:ppl-vs-task}
\end{table}

The automated search across 27 valid configurations reveals two distinct regimes:
\begin{enumerate}
    \item \textbf{Quantization-only} (no eviction): PPL and NIAH agree. Higher quantization increases PPL but NIAH remains 100\%. Both metrics correctly indicate quality.
    \item \textbf{Eviction-included}: PPL degrades gradually ($+$0.97\% to $+$6.30\%) while NIAH exhibits a \emph{binary cliff} --- every eviction configuration produces exactly 60\% NIAH regardless of eviction rate or quantization level. PPL cannot distinguish between 50\% eviction (mild compression) and 70\% eviction (aggressive compression) in terms of retrieval impact.
\end{enumerate}

Critically, both H\textsubscript{2}O (attention-aware) and StreamingLLM (position-based) produce identical 60\% NIAH at 50\% eviction, confirming that the metric gap is \emph{not} an artifact of a poor eviction strategy --- it is an inherent property of the PPL metric itself.

\subsection{Formalization: Task-Aware Distortion}
\label{sec:task-distortion}

We formalize the metric gap using the rate-distortion framework from Section~\ref{sec:info-theory}.

\begin{definition}[Task-Aware Distortion]
For a task with relevant token set $\mathcal{T} \subseteq [N]$ and importance weights $\{w_i\}_{i \in \mathcal{T}}$:
\begin{equation}
    D_{\text{task}}(o, \hat{o}) = \sum_{i \in \mathcal{T}} w_i \|o_i - \hat{o}_i\|^2
\end{equation}
\end{definition}

Since $|\mathcal{T}| \ll N$ for retrieval and QA tasks, $D_{\text{task}} \leq D_{\text{PPL}}$ for any compression scheme. By the monotonicity of the rate-distortion function:
\begin{equation}
    R(D_{\text{task}}) \leq R(D_{\text{PPL}}) \quad \Rightarrow \quad \rho_{\text{task}} \geq \rho_{\text{PPL}}
\end{equation}

The compression ratio achievable under task evaluation is strictly higher than under PPL evaluation. The gap depends on the \emph{information density} of the task:

\begin{equation}
    \frac{\rho_{\text{task}}}{\rho_{\text{PPL}}} \approx \frac{N}{|\mathcal{T}|} \cdot \frac{1}{\max_i w_i}
\end{equation}

For NIAH (1 needle token in $N = 16K$ context): $|\mathcal{T}|/N \approx 1/16000$, but with attention-aware eviction the needle is almost certainly retained, so the effective task distortion is near zero regardless of eviction rate.

\subsection{Implications for the Field}

\paragraph{1. PPL should not be the sole evaluation metric for eviction-based methods.}
PPL is appropriate for evaluating quantization (which affects all tokens uniformly) but is misleading for eviction (which selectively removes tokens). A method that achieves excellent task performance but poor PPL may be strictly preferable for deployment.

\paragraph{2. The ``compression ratio'' of a method depends on the evaluation metric.}
Reporting a single compression ratio is incomplete. We advocate reporting compression ratios at matched quality under both PPL and task-based evaluation, analogous to reporting both PSNR and VMAF in video compression.

\paragraph{3. New evaluation protocols are needed.}
The KV cache compression community would benefit from a standardized evaluation suite that includes:
\begin{itemize}
    \item PPL on standard datasets (for comparability with prior work)
    \item NIAH at multiple context lengths (for retrieval capability)
    \item LongBench/RULER (for diverse long-context tasks)
    \item Per-task compression--quality Pareto curves (not single operating points)
\end{itemize}

\section{Experiments}
\label{sec:experiments}

We evaluate our two-knife compression pipeline (K6V4 asymmetric quantization + H2O attention-aware eviction) across three dimensions: perplexity, retrieval accuracy, and long-context task performance.

\subsection{Setup}

\paragraph{Models.} Llama-3.1-8B-Instruct (GQA, $d=128$, 32 layers), Qwen2.5-7B-Instruct, Ministral-8B-Instruct.

\paragraph{Baselines.}
\begin{itemize}
    \item FP16 (no compression)
    \item TQKV\_6 (TurboQuant uniform 6-bit K+V)
    \item K6V4 (asymmetric: K=6-bit, V=4-bit)
    \item StreamingLLM (position-based eviction)
    \item H2O (attention-aware eviction)
    \item KIVI (INT2 K+V)
\end{itemize}

\paragraph{Metrics.}
\begin{itemize}
    \item \textbf{PPL}: WikiText-2 perplexity (per-token, geometric mean)
    \item \textbf{NIAH}: Needle-in-a-Haystack retrieval accuracy at 4K/8K/16K context
    \item \textbf{LongBench}: Multi-task long-context benchmark (QA, summarization, few-shot)
\end{itemize}

\subsection{Experiment 1: Quantization Baseline}

\begin{table}[h]
\centering
\begin{tabular}{lccccc}
\toprule
\textbf{Model} & \textbf{Config} & \textbf{PPL} & \textbf{$\Delta$PPL} & \textbf{NIAH 4K} & \textbf{Compression} \\
\midrule
\multirow{3}{*}{Llama-3.1-8B} & FP16 & 8.8984 & --- & 5/5 & 1x \\
 & TQKV\_6 & 8.9043 & +0.07\% & 5/5 & 2.67x \\
 & K6V4 & TBD & +0.52\% & TBD & 3.2x \\
\midrule
\multirow{2}{*}{Qwen2.5-7B} & FP16 & TBD & --- & TBD & 1x \\
 & K6V4 & TBD & TBD & TBD & 3.2x \\
\midrule
\multirow{2}{*}{Ministral-8B} & FP16 & 9.3338 & --- & 5/5 & 1x \\
 & TQKV\_6 & 9.3508 & +0.18\% & 5/5 & 2.67x \\
\bottomrule
\end{tabular}
\caption{Quantization-only results. K6V4 achieves 3.2$\times$ compression with $<$1\% PPL increase across all models.}
\label{tab:quant-baseline}
\end{table}

\subsection{Experiment 2: Eviction Gradient (PPL)}

\begin{table}[h]
\centering
\begin{tabular}{lcccc}
\toprule
\textbf{Eviction \%} & \textbf{StreamingLLM (512)} & \textbf{StreamingLLM (16K)} & \textbf{H2O (512)} & \textbf{H2O (16K)} \\
\midrule
0\% & +0.00\% & +0.00\% & +0.00\% & +0.00\% \\
30\% & TBD & $-$0.17\% & TBD & TBD \\
50\% & +0.46\% & $-$0.09\% & TBD & TBD \\
60\% & TBD & $-$0.003\% & TBD & TBD \\
70\% & TBD & +1.04\% & TBD & TBD \\
80\% & TBD & +3.09\% & TBD & TBD \\
85\% & TBD & +6.37\% & TBD & TBD \\
90\% & TBD & TBD & TBD & TBD \\
\bottomrule
\end{tabular}
\caption{PPL increase (\%) at various eviction rates. StreamingLLM data measured; H2O pending. The 1\% PPL cliff shifts right with context length (log scaling, §\ref{sec:ppl-gap}).}
\label{tab:eviction-gradient}
\end{table}

\subsection{Experiment 3: NIAH Retrieval Matrix}

\begin{table}[h]
\centering
\begin{tabular}{lccccc}
\toprule
\textbf{Method} & \textbf{Eviction} & \textbf{Compression} & \textbf{NIAH 4K} & \textbf{NIAH 8K} & \textbf{NIAH 16K} \\
\midrule
FP16 & 0\% & 1x & 100\% & 100\% & 100\% \\
K6V4 only & 0\% & 3.2x & 100\% & TBD & TBD \\
\midrule
K6V4 + StreamingLLM & 50\% & 6.4x & TBD & TBD & 60\% \\
K6V4 + StreamingLLM & 70\% & 9.7x & TBD & TBD & 40\% \\
K6V4 + StreamingLLM & 85\% & 21x & TBD & TBD & 20\% \\
\midrule
K6V4 + H2O & 50\% & 6.4x & TBD & TBD & TBD \\
K6V4 + H2O & 67\% & 9.7x & TBD & TBD & TBD \\
K6V4 + H2O & 75\% & 12.8x & TBD & TBD & TBD \\
K6V4 + H2O & 85\% & 21x & TBD & TBD & \textbf{TBD (key)} \\
K6V4 + H2O & 90\% & 32x & TBD & TBD & TBD \\
\bottomrule
\end{tabular}
\caption{NIAH retrieval accuracy at various compression levels. The key experiment is H2O at 85\% eviction, 16K context --- if NIAH remains $>$80\%, the PPL--task gap is confirmed.}
\label{tab:niah-matrix}
\end{table}

\subsection{Experiment 4: LongBench Task Performance}

\begin{table}[h]
\centering
\begin{tabular}{lcccccc}
\toprule
\textbf{Method} & \textbf{$\rho$} & \textbf{Avg} & \textbf{SingleQA} & \textbf{MultiQA} & \textbf{Summ.} & \textbf{FewShot} \\
\midrule
FP16 & 1x & TBD & TBD & TBD & TBD & TBD \\
K6V4 & 3.2x & TBD & TBD & TBD & TBD & TBD \\
K6V4 + H2O 50\% & 6.4x & TBD & TBD & TBD & TBD & TBD \\
K6V4 + H2O 75\% & 12.8x & TBD & TBD & TBD & TBD & TBD \\
K6V4 + H2O 85\% & 21x & TBD & TBD & TBD & TBD & TBD \\
K6V4 + H2O 90\% & 32x & TBD & TBD & TBD & TBD & TBD \\
\bottomrule
\end{tabular}
\caption{LongBench scores at various compression levels. Task-based evaluation should show higher tolerance to eviction than PPL.}
\label{tab:longbench}
\end{table}

\subsection{Experiment 5: The Metric Gap}

\begin{table}[h]
\centering
\begin{tabular}{lcccl}
\toprule
\textbf{Compression} & \textbf{PPL $\Delta$} & \textbf{NIAH} & \textbf{LongBench $\Delta$} & \textbf{Verdict} \\
\midrule
3.2x (quant only) & +0.52\% & 100\% & TBD & All metrics agree \\
6.4x (+ 50\% evict) & TBD & TBD & TBD & TBD \\
12.8x (+ 75\% evict) & TBD & TBD & TBD & TBD \\
21x (+ 85\% evict) & TBD & TBD & TBD & \textbf{Gap confirmation} \\
32x (+ 90\% evict) & TBD & TBD & TBD & TBD \\
\bottomrule
\end{tabular}
\caption{Summary table comparing PPL vs.\ task metrics at each compression level. The ``gap confirmation'' row is the key result: if PPL degrades significantly but NIAH/LongBench hold, the metric gap is established.}
\label{tab:metric-gap-summary}
\end{table}

\subsection{Experiment 6: Context Length Scaling}

\begin{table}[h]
\centering
\begin{tabular}{lcccc}
\toprule
\textbf{Context $n$} & \textbf{PPL 1\% cliff} & \textbf{Max $\rho$ (PPL)} & \textbf{Max $\rho$ (NIAH)} & \textbf{Predicted cliff} \\
\midrule
512 & 53\% & 6.7x & TBD & 53\% (fit) \\
4,096 & TBD & TBD & TBD & 61\% \\
8,192 & TBD & TBD & TBD & 64\% \\
16,384 & 67\% & 9.7x & TBD & 67\% (fit) \\
32,768 & TBD & TBD & TBD & 70\% \\
\bottomrule
\end{tabular}
\caption{Context length scaling validation. Predicted cliff from log model $f(n) = 0.277 + 0.0405 \ln(n)$. Longer contexts allow more aggressive eviction.}
\label{tab:context-scaling}
\end{table}

\subsection{Experiment 7: Attention-Aware vs Position-Based Eviction}

\begin{table}[h]
\centering
\begin{tabular}{lcccc}
\toprule
\textbf{Method} & \textbf{Eviction} & \textbf{PPL $\Delta$} & \textbf{NIAH 16K} & \textbf{Explanation} \\
\midrule
StreamingLLM & 50\% & $-$0.09\% & 60\% & Position-based: needles at 10\%/25\% evicted \\
StreamingLLM & 70\% & +1.04\% & 40\% & Position-based: needles at 10\%/25\%/50\% evicted \\
StreamingLLM & 85\% & +6.37\% & 20\% & Only 90\% position needle survives \\
H2O & 50\% & TBD & TBD & Attention-based: retains needle \\
H2O & 85\% & TBD & TBD & Attention-based: retains needle \\
\bottomrule
\end{tabular}
\caption{Head-to-head comparison of eviction strategies. StreamingLLM optimizes PPL; H2O optimizes task performance. Both may achieve similar PPL, but H2O should dramatically outperform on NIAH.}
\label{tab:eviction-comparison}
\end{table}
  % Tables with TBD placeholders
% \section{Conclusion}
\label{sec:conclusion}

We have presented a systematic empirical study of KV cache compression across five models spanning 7B--70B parameters, three GQA ratios, and two compression paradigms (quantization vs.\ eviction). Our findings revise several assumptions in the KV cache compression literature:

\paragraph{1. Quantization safety is not universal.}
While 4-bit KV quantization is NIAH-safe for GQA 4:1 models (Llama, Mistral: $\leq$3\% PPL increase, 100\% NIAH), it catastrophically fails on GQA 7:1 models (Qwen: PPL $\to$6615, NIAH $\to$0\%). The root cause is K-quantization: V tolerates 4-bit universally, but K errors enter the softmax nonlinearity where they are amplified by all shared query heads. We propose a \textbf{GQA-aware quantization rule}: K precision must increase with GQA ratio, while V precision can remain aggressive.

\paragraph{2. Eviction has a fundamental causal limitation.}
Both attention-aware (H\textsubscript{2}O) and position-based (StreamingLLM) eviction produce identical NIAH degradation at matched eviction rates. This is not an algorithmic limitation---it is a causal one. Eviction decisions are made at prefill time without access to future queries. We survey 20+ methods for predicting future attention (\S\ref{sec:theory}) and find that key vectors carry only 0.12 bits of predictive information about future attention patterns~\citep{limits-scoring-2026}, establishing the difficulty as information-theoretic.

\paragraph{3. Perplexity is a misleading metric for eviction.}
At 50\% eviction, PPL improves ($-$0.09\%) while NIAH drops from 100\% to 60\%---a dangerous false positive. Quantization produces the opposite signature: PPL degrades (+3\%) while NIAH holds at 100\%. These orthogonal responses mean neither metric alone characterizes compression quality. We recommend \textbf{dual-metric reporting} (PPL + retrieval) as the minimum standard.

\paragraph{4. Three orthogonal compression axes.}
KV cache compression operates along three independent axes: \emph{dimension reduction} (low-rank projection, 1.5--2$\times$), \emph{bit reduction} (quantization, 2.67--4$\times$), and \emph{token reduction} (eviction, 2--5$\times$). No prior work systematically studies their interaction. Our experiments show the axes are largely orthogonal, suggesting multiplicative compression is achievable---but the token reduction axis is constrained by the causal limitation for retrieval-critical applications.

\subsection{Practical Recommendations}

Based on our cross-model empirical results:

\begin{enumerate}
    \item \textbf{Always use asymmetric K/V quantization.} K is more sensitive than V due to softmax amplification. Set K precision $\geq$ V precision, with the gap increasing for higher GQA ratios.
    \item \textbf{Profile outlier layers before deployment.} A single outlier layer (e.g., Qwen Layer 0, $K_\text{max}=93$) can render quantization catastrophic. Skip these layers or use higher precision.
    \item \textbf{Do not use PPL alone to evaluate eviction.} PPL can show improvement while retrieval collapses. Always include a retrieval task (NIAH or equivalent).
    \item \textbf{Match eviction aggressiveness to context length.} The 1\% PPL cliff shifts from $\sim$53\% at 4K to $\sim$67\% at 16K, following $f(n) = 0.277 + 0.0405 \ln(n)$.
    \item \textbf{Do not assume cross-model consistency for eviction.} Mistral tolerates 85\% eviction where Llama fails at 67\%. Quantization behavior is more consistent ($\sim$3\% at q4\_0 across 8B--70B).
\end{enumerate}

\subsection{Limitations}

Our study has several limitations. First, all experiments use llama.cpp's GGUF quantization types; results may differ with calibrated quantization methods (GPTQ, AWQ) that can handle outliers more gracefully. Second, our NIAH evaluation uses a single-needle synthetic benchmark; multi-needle or natural retrieval tasks may show different sensitivity profiles. Third, we do not evaluate methods requiring fine-tuning (DMC, LookaheadKV), which may achieve better compression--quality trade-offs at the cost of training compute. Fourth, eviction experiments use StreamingLLM and H\textsubscript{2}O only; more sophisticated eviction methods (SnapKV, PyramidKV) may shift the cliff positions.

\subsection{Future Work}

Several directions emerge from our findings:
\begin{itemize}
    \item \textbf{Selective recompute:} Store all tokens at ultra-low precision plus token IDs; at decode time, use low-precision attention to identify top-$k$ important tokens and recompute their full-precision KV on the fly. Our theoretical analysis suggests 34.6$\times$ compression at 100\% NIAH with $\sim$11\% compute overhead (\S\ref{sec:theory}).
    \item \textbf{Expected Attention for NIAH-safe eviction:} Model future queries as Gaussian-distributed and compute closed-form expected attention scores---a training-free approach to the causal limitation.
    \item \textbf{GQA-aware compression toolkit:} An open-source tool that automatically profiles model architecture, detects outlier layers, and recommends optimal asymmetric K/V configurations.
    \item \textbf{PALU-style low-rank stacking:} Adding a fused linear dimension reduction layer (1.5--2$\times$) on top of quantization and eviction as a fourth compression axis.
\end{itemize}
   % TODO

\bibliographystyle{plainnat}
% \bibliography{references}

\end{document}
